\subsection{Two far endpoints at inverse-logarithmic margin}
\label{subsec:random-norm-tags}

A logarithmic-size selection of norm fibers gives the inverse-logarithmic
capacity margin of the quotient construction while retaining almost every
native challenge.  The evaluation domain is chosen for the construction.
Crites--Stewart~\cite[Corollary~1]{crites-stewart2025} already give
every affine challenge nearby at this order of margin, over arbitrary
fields and domains, with a far direction. The point of the following
specialization is its two far endpoints and exact agreement profile;
no novelty claim for that refinement is made.

% Standalone input after Lemma lem:fixed-cardinality-character-products.
% Natural logarithms throughout.  This proves existence for one fixed pole b.
\begin{lemma}\label{lem:random-tag-character-bias}
Let $E=\mathbb F_q$, where $q=p^d$ and $d\ge2$, and let
$b\in E$ have degree $d$ over $\mathbb F_p$.
Fix a set $T\subset\mathbb F_p$ of $e$ excluded tags, put
$U=\mathbb F_p\setminus T$, and define
\[
 N=p-e,\qquad \mu=\frac{(d-1)\sqrt p+e}{p-e}.
\]
For $1\le s\le N$, choose an $s$-element subset $G$ of $U$ uniformly
without replacement.  For any real $\beta>\mu$,
\[
 \Pr\left[\exists\chi\ne1:
   \left|\sum_{a\in G}\chi(b-a)\right|>\beta s\right]
 \le 4(q-2)\exp\left(-\frac{s(\beta-\mu)^2}{4}\right),
\]
where the quantifier ranges over all nontrivial multiplicative characters
of $E^*$.
\end{lemma}
\begin{proof}
Katz's affine-line estimate~\cite[Theorem~1]{katz1989-character-sums}
and deletion of the $e$ tags give
$|N^{-1}\sum_{a\in U}\chi(b-a)|\le\mu$ for every nontrivial $\chi$.
We prove the needed concentration estimate without assuming independence
of the sampled tags.

For a real population $x_1,\ldots,x_N\in[-1,1]$ with mean $\bar x$,
let $S$ be the sum over a uniform $s$-element subset.  For real $t$,
\[
 \mathbb E e^{tS}
 =\frac{e_s(e^{tx_1},\ldots,e^{tx_N})}{\binom Ns}
 \le\left(\frac1N\sum_{i=1}^N e^{tx_i}\right)^s.
\]
Here $e_s$ denotes the elementary symmetric polynomial.  Its displayed
bound follows by averaging pairs of positive arguments with fixed sum:
with all other arguments fixed, it has the form
$A y_i y_j+B(y_i+y_j)+D$ with $A\ge0$.
Equivalently, its maximum on the nonnegative simplex occurs when all
arguments are equal.  Indeed, for $s\ge2$ a maximizing tuple has at least
$s$ positive arguments, so averaging any unequal pair strictly increases
the value; the case $s=1$ is immediate.

If $X$ is one uniform population entry, the second derivative of
$\log\mathbb E e^{t(X-\bar x)}$ is the variance under an exponential
tilt, and is at most $1$ because $X\in[-1,1]$.  Its value and first
derivative vanish at $t=0$.  Thus
$\mathbb E e^{t(S-s\bar x)}\le e^{st^2/2}$.
Chernoff's bound, applied with both signs of $t$, yields
\[
 \Pr[|S-s\bar x|>su]\le2e^{-su^2/2}.
\]
Apply this separately to the real and imaginary parts of
$\chi(b-a)$.  A complex deviation of modulus greater than $sv$ has
a coordinate of absolute value greater than $sv/\sqrt2$, so its
probability is at most $4e^{-sv^2/4}$.
Set $v=\beta-\mu$ and take the union bound over the $q-2$ nontrivial
characters.
\end{proof}

\begin{corollary}\label{cor:logarithmic-tag-product-coverage}
Fix $0<\eta\le1/2$, set $\kappa=\eta(1-\eta)$ and
\[
 C=\max\{128,8/\kappa\},\qquad s=\lceil C\log q\rceil.
\]
Fix $a_0\in\mathbb F_p^*$, and suppose
\[
 q\ge C+2,\qquad s\le p-2,\qquad
 \frac{(d-1)\sqrt p+2}{p-2}\le\frac14.
\]
With probability at least $1-4/q$, a uniform $s$-element subset
$G\subset\mathbb F_p^*\setminus\{a_0\}$ has the following property
simultaneously for every integer $r$ with $\eta\le r/s\le1-\eta$:
every element of $E^*$ is a product of exactly $r$ distinct elements of
$\{b-a:a\in G\}$.
\end{corollary}
\begin{proof}
Use Lemma~\ref{lem:random-tag-character-bias} with $e=2$ and
$\beta=1/2$.  The failure probability is at most
\[
 4(q-2)e^{-s/64}\le4q^{1-C/64}\le4/q.
\]
For any successful $G$, Lemma~\ref{lem:fixed-cardinality-character-products}
applies with $B=s/2$.  Its total nontrivial-character error is at most
\[
 (q-2)(s+1)e^{-\kappa s/2}
 \le\frac{s+1}{q^3}\le\frac1q<1.
\]
For the last inequality, use $s+1\le C\log q+2\le Cq+2\le q^2$,
which follows from $q\ge C+2$.
One simultaneous character-bias event suffices for every indicated $r$;
no further probability union bound is needed.
\end{proof}


\begin{theorem}\label{thm:random-norm-tag-log-gap}
Fix $0<\rho<1$, an integer $d\ge2$, and
\[
 C>\max\left\{64,\frac4{\rho(1-\rho)}\right\}.
\]
For every sufficiently large prime $p$, put
\[
 q=p^d,\quad E=\mathbb F_q,\quad
 m=\frac{q-1}{p-1},\quad s=\lceil C\log q\rceil,\quad
 n=(s+1)m,\quad J=\lfloor\rho n\rfloor.
\]
There is an $n$-point domain $D\subset E$ and words $f,g:D\to E$ with
\[
 \operatorname{agr}_J(f)=\operatorname{agr}_J(g)
 =\operatorname{CA}_J(f,g)=J+m-1,
\]
such that every $\lambda\in E\setminus\{0\}$ satisfies
\[
 \operatorname{agr}_J(f+\lambda g)\ge J+2m-1.
\]
The capacity margin and source-agreement gap are respectively
\[
 \frac{2m-1}{n}=\frac2{s+1}-\frac1n,
 \qquad \frac mn=\frac1{s+1},
\]
both of order $1/\log n$.  The alphabet and characteristic obey
\[
 q=\Theta\!\left((n/\log n)^{d/(d-1)}\right),\qquad
 p=\Theta\!\left((n/\log n)^{1/(d-1)}\right).
\]
\end{theorem}

\begin{proof}
Choose $b\in E$ of degree $d$ over $\mathbb F_p$ and fix
$a_0\in\mathbb F_p^*$.  For every nontrivial multiplicative character
$\chi$ of $E^*$, Katz's estimate~\cite[Theorem~1]{katz1989-character-sums}
gives
\[
 \left|\frac1{p-2}
 \sum_{a\in\mathbb F_p^*\setminus\{a_0\}}\chi(b-a)\right|
 \le\frac{(d-1)\sqrt p+2}{p-2}\le\frac14
\]
for sufficiently large $p$.
Choose an $s$-element subset $G$ of this population uniformly without
replacement.  Lemma~\ref{lem:random-tag-character-bias} and a union bound
show that, with probability at least
$1-4(q-2)\exp(-s/64)>0$, it satisfies
\[
 \left|\sum_{a\in G}\chi(b-a)\right|\le s/2
 \quad\hbox{for every nontrivial $\chi$}.
\]
Fix such a $G$, write $N(X)=X^m$, and set
$D=N^{-1}(G\cup\{a_0\})$.  Each nonzero norm fiber has size $m$,
so $|D|=(s+1)m=n$.

Write $J-1=(r-2)m+w$, where $0\le w<m$, and put $\theta=r/s$.
Then $\theta\to\rho$ and $2\le r<s$ for sufficiently large $p$.
Lemma~\ref{lem:fixed-cardinality-character-products}, with
$\mathcal A=\{b-a:a\in G\}$ and $B=s/2$, shows that all elements of
$E^*$ are products of exactly $r$ distinct elements of $\mathcal A$,
provided
\[
 (q-2)(s+1)\exp\bigl(-\theta(1-\theta)s/2\bigr)<1.
\]
Our choice of $C$ makes both displayed probability conditions hold
eventually.  Also $s<p-2$ eventually, so the sampling is defined.

Choose a $w$-point set $B_0\subset N^{-1}(a_0)$, put $R=L_{B_0}$,
and use the norm quotient words
\[
 f=R\frac{N^r-b^r}{N-b},\qquad g=-\frac R{N-b}
 \quad\hbox{on }D.
\]
The denominator is nonzero on $E$.  The polynomial $f$ is monic of
degree $J+m-1$, and the agreement equation for $g$ is the nonzero
polynomial $R+(N-b)h=0$, of degree at most $J+m-1$ for every
$\deg h<J$.  These source bounds are attained simultaneously by
interpolating the two quotient-variable functions on any $r-1$ tags
of $G$, substituting $N$, and multiplying by $R$.  The resulting
degree-at-most-$J-1$ witnesses match on those fibers and $B_0$,
giving exactly $J+m-1$ common agreements.

For $\lambda\ne0$, choose $S\subset G$, $|S|=r$, with
$V_S(b)=-\lambda$, where $V_S(Y)=\prod_{a\in S}(Y-a)$.
Writing $P_S(Y)=Y^r-V_S(Y)$, the polynomial
\[
 h_S=R\frac{P_S(N)-P_S(b)}{N-b}
\]
has degree at most $w+(r-2)m=J-1$ and satisfies
\[
 f+\lambda g-h_S=\frac{R V_S(N)}{N-b}.
\]
Its zeros on $D$ are exactly $B_0$ and the $r$ norm fibers indexed
by $S$, totaling $w+rm=J+2m-1$.
Finally, $n=\Theta(p^{d-1}\log p)$ proves the parameter assertions.
\end{proof}

For finite parameters the proof requires only $s<p-2$, $2\le r<s$,
and the three explicit inequalities
\[
 \frac{(d-1)\sqrt p+2}{p-2}\le\frac14,\qquad
 4(q-2)e^{-s/64}<1,\qquad
 (q-2)(s+1)e^{-\theta(1-\theta)s/2}<1.
\]
At threshold $T=J+2m-1$, exactly $q-2$ of the affine mixtures
$(1-t)f+tg$, $t\in E$, are nearby; the two endpoints are farther by
$m/n$.  This gives almost-complete coverage at an inverse-logarithmic
margin by standard quotient and character-sum ingredients.
The domain is a chosen union of norm fibers, without a subgroup
requirement.  The alphabet remains an extension field, and the domain
length is $\Theta(p^{d-1}\log p)$, rather than $p^d$.
No prescribed short evaluation domain is obtained.
